\section{Indledning}

% alt + z for line wrap :)

Problemfelt

Den 23. oktober 1944 trængte Gestapo ind på Viborg Katedralskole og arresterede 14 modstandsfolk, mens flere andre måtte flygte(\cite{katedralskole2026}). Netop sådanne begivenheder kan gøre historien nærværende for elever, men erfaringer fra undervisningsverdenen tyder på, at potentialet sjældent forløses. En debattør med 40 års erfaring fra lærerseminariet beskrev i 2017 historieundervisningen som præget af "en gammel tradition om at repetere nationalt vedtagne fortællinger om fædrelandet" og påpegede, at "børnene får ikke lov at mærke historien; den skal bare læres" (\cite{ronhede2017historiefaget}). Selvom der er tale om et debatindlæg, stemmer kritikken overens med international forskning. Elever engageres især, når historieundervisningen gøres relevant for deres egen livsverden, og når den inddrager kontekstualisering, argumentation og samarbejde(\cite{VanDoorsselaere2025}). Samtidig viser en voksende mængde studier, at spilbaserede oplevelser kan øge både motivation og historieforståelse markant. Et studie af det AR-understøttede spil Exploring Ancient Greece fandt, at elever i alderen 10-14 år forbedrede deres historieviden signifikant og vurderede spillet som en foretrukken læringsform(\cite{Bekas2024}). En metaanalyse dækkende 39 studier konkluderede desuden, at digitale læringsspil generelt er mere motiverende end traditionel undervisning, selvom effekten på faktuel videnstilegnelse er blandet(\cite{Wouters2013}). Projektet samarbejder med Viborg Katedralskole. Skolen skal i forbindelse med bygningens 100-års jubilæum fejre noget. Gennem samtaler med undervisere har projektgruppen erfaret, at der på skolen er begejstring for et digitalt spil, der inddrager eleverne aktivt og skaber en fælles oplevelse omkring skolens historie. \\ En central del af problemfeltet er at forstå, hvornår et spil er den rigtige ramme for historieformidling. Mayer, Makransky og Terkildsen har påpeget, at der stadig mangler empirisk evidens for, at øget teknologisk immersion i sig selv skaber bedre læring(\cite{Mayer2022}) . I deres metaundersøgelse rapporterede studerende højere tilstedeværelse i immersive VR-betingelser, men lærte mindre end i mindre immersive formater. Her adskiller projektets ambition sig: målet er ikke primært faktuel indlæring. En nyere analyse af 118 publikationer underbygger, at spilbaseret læring har et markant potentiale for at styrke elevers engagement og historiske forståelse, forudsat at den didaktiske ramme omkring spillet er velovervejet (\cite{Lai2025}).


Projektets problemformulering er derfor:

Hvordan kan man hylde fejringen af Katedralskolens 100-års jubilæum for bygningen ved hjælp af et spil?
Hertil knytter sig følgende underspørgsmål:
\begin{itemize}
    \item Hvordan kan vi fremvise bygningen i et historisk perspektiv?
    \item Hvordan kan vi inddrage og samarbejde med nuværende elever på katedralskolen i forbindelse med detaljer og kreative tilføjelser for at engagere elever til at lære om skolens historik?
    \item Hvordan kan vi designe en digital oplevelse med udgangspunkt i en historisk begivenhed?
\end{itemize}
